\newpage
%%%%%%%%%%%%%%%%%%%%%%%%%%%%%%%%%%%%%%%%%%%%%%%%%%
%%%%%%%%%%%%%%%%%%%%%%%%%%%%%%%%%%%%%%%%%%%%%%%%%%
\section{Using only a sequence to zero} \label{sec:nullseq}
%%%%%%%%%%%%%%%%%%%%%%%%%%%%%%%%%%%%%%%%%%%%%%%%%%
%%%%%%%%%%%%%%%%%%%%%%%%%%%%%%%%%%%%%%%%%%%%%%%%%%

Looking back, the Archimedean property (V3) of $\Q$ was used in only two
places: to produce the null sequence $(e_n)$ (Lemma~\ref{lem:nullseq}), and to
prove (V3) for $\Qb$ (Step 3c). Everything else used only the existence of
$(e_n)$. So (V3) may be replaced by the weaker assumption
\begin{enumerate}[nosep,leftmargin=4em]
    \item[(V3$'$)] there is a sequence $e_n>0$ in $V$ with
    $\forall\epsilon>0\ \exists N\ \forall n\ge N,\ e_n<\epsilon$;
\end{enumerate}
such a $V$ is automatically dense (there is no smallest positive element),
and $(e_n)$ may be chosen with $e_{n+1}+e_{n+1}\le e_n$ (replace $e_{n+1}$ by
$\min(e_{n+1},e_n/2)$).

\begin{theorem} (completion without Archimedes) $\mid$ \label{thm:V3prime}
Let $\Q$ be nontrivial and satisfy (V1), (V2) and (V3$'$). Then
Sections~\ref{sec:construct}--\ref{sec:smallest} hold verbatim, except that
$\Qb$ satisfies (V3$'$), with null sequence $\iota(e_n)$
(Corollary~\ref{cor:wise}), in place of (V3).
\end{theorem}

The completion need no longer be $\R$:

\begin{example} ($\Q\times\Q$ completes to $\Q\times\R$) $\mid$ \label{ex:lex}
Let $V=\Q\times\Q$ (now $\Q$ the rationals) with componentwise addition and the
\emph{lexicographic} order: $(a,b)>0$ iff $a>0$, or $a=0$ and $b>0$. It
satisfies (V1), (V2) and (V3$'$) with $e_n=(0,(\tfrac12)^n)$, but not (V3):
$n\cdot(0,1)=(0,n)<(1,0)$ for every $n$. For a Cauchy sequence $(a_n,b_n)$,
taking $\epsilon=(0,1)$ shows that $a_n$ is eventually constant, say $=a$, and
taking $\epsilon=(0,\delta)$ shows that $(b_n)$ is Cauchy in $\Q$. Hence
$[(a_n,b_n)]\mapsto(a,[(b_n)])$ identifies $\Qb$ with $\Q\times\R$
(lexicographic), which is not Archimedean.
\end{example}


\section{Completing any metric space}

Say I have some arbitrary metric $(X,d: X \times X \to A)$ for some well-valued A. We wish to find the smallest metric that completes this one. Under the strength of the contents of this paper hitherto, the completion of A is either isomorphic to $\Z$ (th integers) or to $\R$ (the reals, $\Q$). However the integers isomorphically map into the reals so there is no harm in just always choosing the resulting completion metric to take values in $\R$....\\
Our goal is, given a metric $(X,d_X: X \times X \to \R)$, find a larger space $(\bar{X},d_{\bar{X}}:\bar{X} \times \bar{X} \to \R)$ that (1) is complete (2) contains isometrically $(X,d)$, and (3) is the smallest set that does (1) \& (2).  \\

The result is immediate as the proofs are exactly the same. We define 

\[ \bar{X} = \{ [(a_n)] \mid\; \text{ $(a_n)$ is cauchy sequence in $(X,d_X)$}\} \]
\[ d_{\bar{X}} ([(a_n)], [(b_n)]) = [(d_X(a_n,b_n))_{n \in \N}] = \lim_{n \ \in \N, n \to \infty}d_X(a_n,b_n) \]

The equivalence relation required in eq \_\_ is also the same as is required by our distance function needing to be a metric. 
Because nothing in our proofs of completeness used any of the structure on the metric space $\Q$ (only used this stuff on the values space), the same proofs follow through. 

Note the reason for denoting $\bar{X}$ and $\bar{Q}$ is because these sets are the closure of the corresponding $X$ and $Q$ in their respective metric spaces. 


\section{What sets complete to R?}

add multiplication to R. Everything follows through from added multiplication from $Q$ and we define 
\[ XY := ... \]
The isormphism $l$ also preserves this multiplication so it all works out. Now we have a notion of the number $1$. Such a set $\bar{Q},+,<,*,\bar{d}(X,Y) := |X - Y| $ is ussually denoted by the familiar $\R$. \\

example sets: 

\[ A := \{ \frac{z}{2^n} \mid z \in \Z, n \in \N \} \]


Given $X \in \R, X \geq 0$ By archemedian, given an n, 

\[ \exists N \mid 2^nX < N \]

Explicitely choose the minimal such N and denote $N_{0,(n)} := N_{\text{minimal}} - 1$. We have, 

\[ N_{0,(n)} := \lfloor 2^n X\rfloor \leq 2^nX < N_{0,(n)} + 1 \]

We get the resulting sequence $\frac{\lfloor 2^n X\rfloor}{2^n}$ which is a cauchy sequence into $A$ and whose limit / generating equivalence class gives exactly X. This gives an explicit construction of a element of a real number without using choice, as promised. \\


decimal representation, base anything representation, etc. all just different versions of finding explicit formulas for the equivalence class of Cauchy sequences that makes a real number a real number. \\